% ====================================================================
% performance.tex (design note example)
% ====================================================================

\section{Predicted Performance}
\label{sec:performance}

\subsection{Frequency response}

Substituting $s = \mathrm{j}2\pi f$ into \cref{eq:transferFunction} gives the
magnitude of the response at any frequency,

\begin{equation}
  \label{eq:magnitude}
  \lvert H(f) \rvert =
  \frac{1}{\sqrt{\bigl(1 - (f/f_0)^2\bigr)^2 + \bigl(f/(Q f_0)\bigr)^2}},
\end{equation}
which is plotted in \cref{fig:filterResponse} for the nominal component
values. The response is \qty{0.27}{\decibel} down at the top of the signal
band and \qty{39.2}{\decibel} down at \qty{9.5}{\kilo\hertz}, so both
requirements are met with the nominal values.

% ====================================================================
% filterResponse.tex
% --------------------------------------------------------------------
% Magnitude response of the filter at its nominal values, drawn with
% pgfplots from the design equation (f0 = 995.9 Hz, Q = 0.7071).
% ====================================================================

\begin{figure}[htbp]
\centering
\begin{tikzpicture}
\begin{semilogxaxis}[
  width=0.92\linewidth,
  height=7.5cm,
  grid=both,
  minor grid style={gray!20},
  major grid style={gray!45},
  xmin=10, xmax=20000,
  ymin=-60, ymax=5,
  ytick={-60,-50,-40,-30,-20,-10,0},
  xlabel={Frequency (Hz)},
  ylabel={Gain (dB)},
  legend pos=south west,
  legend cell align=left,
]
  \addplot[domain=10:20000, samples=300, very thick, linkblue]
    {-10*log10((1 - (x/995.9)^2)^2 + (x/(995.9*0.7071))^2)};
  \addlegendentry{Filter response}

  % Stopband requirement: 35 dB at 9.5 kHz
  \addplot[only marks, mark=*, mark size=2.5pt, red] coordinates {(9500,-35)};
  \addlegendentry{Requirement R5}

  % Top of the signal band
  \draw[dashed, gray] (axis cs:500,-60) -- (axis cs:500,5);
  \node[font=\footnotesize, anchor=south east, rotate=90]
    at (axis cs:500,-58) {top of signal band};

  \node[font=\footnotesize, anchor=west]
    at (axis cs:1100,-3) {\qty{-3}{\decibel} at \qty{996}{\hertz}};
  \node[font=\footnotesize, anchor=east]
    at (axis cs:9000,-41) {\qty{-39.2}{\decibel}};
\end{semilogxaxis}
\end{tikzpicture}
\caption{Calculated magnitude response at the nominal component values. The
curve passes \qty{4.2}{\decibel} below the stopband requirement.}
\label{fig:filterResponse}
\end{figure}


\subsection{Effect of component tolerance}

The nominal figures only show that the design is sound on paper, as every
board will be built with parts that are a little off their marked values. The
response was therefore recalculated with each of the four values at either
end of its tolerance band, which gives sixteen combinations, and the worst
result from those is shown against the requirement in \cref{tab:tolerance}.

\begin{table}[htbp]
\centering
\begin{tabularx}{\textwidth}{@{} X l l l @{}}
\toprule
\textbf{Quantity} & \textbf{Requirement} & \textbf{Nominal} & \textbf{Worst case} \\
\midrule
Natural frequency, $f_0$                 & n/a                    & \qty{996}{\hertz}    & \qtyrange{939}{1059}{\hertz} \\
Quality factor, $Q$                      & n/a                    & 0.707                & 0.673 to 0.743 \\
Gain at \qty{500}{\hertz}                & Within \qty{1}{\decibel} & \qty{-0.27}{\decibel} & \qty{-0.49}{\decibel} \\
Peaking in the passband                  & Within \qty{1}{\decibel} & None                 & \qty{0.04}{\decibel} \\
Attenuation at \qty{9.5}{\kilo\hertz}    & \qty{35}{\decibel} or more & \qty{39.2}{\decibel} & \qty{38.1}{\decibel} \\
\bottomrule
\end{tabularx}
\caption{Predicted performance at the nominal values and at the worst
combination of component tolerances.}
\label{tab:tolerance}
\end{table}

The stopband requirement is met with \qty{3.1}{\decibel} in hand in the worst
case, and the passband stays within half a decibel. The natural frequency
moves by about six per cent either way, almost all of which is due to the
capacitors, so tighter resistors would gain very little.

Taking every part to its limit at once is pessimistic, so the calculation was
repeated as a Monte Carlo run of 20,000 builds, with each value drawn at
random from within its tolerance band. The standard deviation of $f_0$ was
\qty{21}{\hertz}, and no build fell outside the worst-case figures in
\cref{tab:tolerance}.

\subsection{Limits of this analysis}

These figures assume an ideal op-amp. In a real Sallen-Key filter the
attenuation does not keep increasing with frequency, because at high
frequencies the op-amp's output impedance rises and the input signal reaches
the output through $C_1$~\cite{karki_sallen_key_2002}. With the part
specified this is expected to matter only well above \qty{10}{\kilo\hertz},
but it cannot be confirmed by calculation alone, so it is one of the things
the tests in \cref{sec:verification} are there to check.
